\chapter{Non-standard physics}\label{ch:nonstandard}

GLoBES computes event rates from oscillation probabilities, and nothing else
in the simulation depends on how these are computed. To simulate a model
beyond three-flavour oscillations (sterile neutrinos, non-standard
interactions, decoherence, decay, \dots) replace the probability engine.

\begin{apifun}{probability\_engine!(P!, extra; setup=current\_setup())\\
standard\_engine!(; setup=current\_setup())}
Use \code{P!} to compute probabilities, with the additional parameters named
in \code{extra}, or go back to the standard engine. \code{P!(P, params, E,
antineutrino, lengths, densities)} fills the $3\times3$ matrix
\code{P[α, β]} $=P(\nu_\alpha\to\nu_\beta)$ for the \code{OscParams}
\code{params} at energy \code{E} (GeV) through the layers of the matter
profile (km and $\mathrm{g/cm^3}$, already scaled by $\rho$).
\end{apifun}

\noindent After the engine is replaced, \code{OscParams} must have the
additional parameters (\code{extra=(name=value, ...)}), and the data must be
simulated again. The additional parameters work everywhere the standard ones
do: they can be fitted, have priors, and be scanned.

As an example, here is an engine that multiplies the standard probabilities
by $1-\epsilon$, a crude model of neutrino decay or of a loss of
normalisation:
\begin{julia}
engine = GLoBES.StandardEngine()
probability_engine!((:ϵ,)) do P, p, E, anti, lengths, densities
    GLoBES.set_oscillation_parameters!(engine, p.values[1:6])
    GLoBES.probability_matrix!(P, engine, anti ? -1 : +1, E, length(lengths),
                               lengths, densities, -1.0)
    P .*= 1 - p.ϵ
end

truth = OscParams(truth; extra=(ϵ=0.0,))
simulate!(truth)
r = marginalize(OscParams(truth; ϵ=0.1); free=(:ϵ, :θ23), priors=(θ23=0.05,))
r.params.ϵ                     # ≈ 0: the data have no loss
\end{julia}
The standard engine used here (\code{GLoBES.StandardEngine},
\code{GLoBES.set\_oscillation\_parameters!} and
\code{GLoBES.probability\_matrix!}) is the one of Chapter~\ref{ch:prob}: a
model that modifies the Hamiltonian can build on its code in
\file{src/probability.jl}.

\begin{note}
The engine is called for every sampling point of every experiment for every
$\chi^2$ evaluation, so it should be fast and should not allocate memory.
\end{note}
